% Additive integration proposal for the existing derived-source mechanism.
% Not loaded by any current reader; do not overwrite the frozen source.
\OLSADeclareDocumentTextCorrection{OLP-0005}{SF-0005-PERFECT-DOMAIN}{عددی را \emph{کامل} می‌نامیم اگر و تنها اگر با مجموع مقسوم‌علیه‌های
سرهٔ خود برابر باشد (یعنی عددهایی که آن را بدون باقیمانده تقسیم می‌کنند
اما با خود آن عدد برابر نیستند). برای نمونه، $6$ کامل است زیرا
مقسوم‌علیه‌های سرهٔ آن $1$، $2$ و~$3$ هستند و $6 = 1 + 2 + 3$.
}{عدد صحیح مثبتی را \emph{کامل} می‌نامیم اگر و تنها اگر با مجموع مقسوم‌علیه‌های
سرهٔ مثبت خود برابر باشد (یعنی عددهای صحیح مثبتی که آن را بدون باقیمانده تقسیم می‌کنند
اما با خود آن عدد برابر نیستند). برای نمونه، $6$ کامل است زیرا
مقسوم‌علیه‌های سرهٔ مثبت آن $1$، $2$ و~$3$ هستند و $6 = 1 + 2 + 3$.
}
\OLSADeclareDocumentTextCorrection{OLP-0005}{SF-0005-EXTENSIONALITY-EXISTENCE}{اصل گسترش تضمین می‌کند که همواره تنها یک مجموعه از $x$هایی وجود دارد
که $\phi(x)$ درباره‌شان صادق است. پس اصل گسترش مجوز می‌دهد
$\Setabs{x}{\phi(x)}$ را \emph{یگانه} مجموعهٔ $x$هایی بنامیم
که~$\phi(x)$ درباره‌شان صادق است.}{اصل گسترش تضمین می‌کند که حداکثر یک مجموعه از $x$هایی وجود دارد
که $\phi(x)$ درباره‌شان صادق است. پس، اگر چنین مجموعه‌ای وجود داشته باشد،
اصل گسترش مجوز می‌دهد $\Setabs{x}{\phi(x)}$ را
\emph{یگانه} مجموعهٔ $x$هایی بنامیم که~$\phi(x)$ درباره‌شان صادق است.\footnote{یادداشت ویراستاری: در مثال عدد کامل، مثبت‌بودن عدد و مقسوم‌علیه‌ها تصریح شده است. همچنین شرط وجود مجموعه در بیان اصل گسترش تصریح شده است؛ این اصل به‌تنهایی وجود مجموعه را تضمین نمی‌کند.}}
\OLSADeclareSetup{OLP-0005}{\olsa_source_correction_install_document_text:}
